% =============================================================================
% Pillar 3 — Iter 131 — G=4 vs G=32 vs the Broader Sweep, in relation to
% the 'It Takes Two: Your GRPO Is Secretly DPO' claim (arXiv 2510.00977)
% that G=2 ~ G=16 with 97.6% retention. We test whether G=4 ~ G=32
% holds at broader scale (T-budget sweep).
% Driver: scripts/group_size_analysis.py
% Inputs:  experiments/results/groupsize_zvf_sweep.json      (G in {2,4,8,16}, n=3)
%          experiments/results/group_size_token_normalized.tsv (G x T grid 5x4, n=20)
% Outputs: experiments/results/group_size_effect.tsv
%          figures/group_size.pdf
% =============================================================================
\section{Group Size: G=4 vs G=32 vs the Broader Sweep}
\label{sec:group-size}

\paragraph{Headline question.}
A reviewer (and the recent theoretical literature) asks: \emph{is GRPO's
group size $G$ a real hyperparameter or is it effectively constant?}
Wu et al.~\cite{wu2025grpodpo} prove that GRPO is algebraically
equivalent to an implicit contrastive objective (DPO), with $G$ affecting
only Monte Carlo variance, and report that \textbf{2-GRPO retains
97.6\% of 16-GRPO performance} while requiring only 12.5\% of the
rollouts and 21\% of the training time. The opposite intuition---that
larger $G$ amortizes wasted (all-correct or all-incorrect) rollouts---is
also in the literature. We bring both perspectives to our data.

\paragraph{Summary of evidence.}
Three sources speak to the question, summarized in
Table~\ref{tab:groupsize-summary}:

\begin{enumerate}
  \item \textbf{Measured small-scale sweep} on Qwen2.5-0.5B /
        arithmetic\_correctness, 40 steps, $G \in \{2, 4, 8, 16\}$ at
        three seeds each
        (\texttt{experiments/results/groupsize\_zvf\_sweep.json}).
        Held-out accuracy traces an inverted-U with its numerical apex
        at $G{=}8$ ($0.990{\pm}0.003$); $G{=}2$ ($0.982{\pm}0.004$) and
        $G{=}16$ ($0.978{\pm}0.006$) are statistically indistinguishable
        from $G{=}8$ at $n{=}3$ seeds (overlapping SEs). On this task
        $G{=}2$ actually \emph{retains 100.3\% of 16-GRPO held-out accuracy}
        (\texttt{experiments/results/group\_size\_effect\_dpo\_check.tsv}),
        comfortably above the 97.6\% Wu et al.\ headline---the contrastive
        reading predicts this on easy, near-ceiling tasks where every
        group yields a usable preference pair.
  \item \textbf{Token-budget normalized sweep} (illustrative reanalysis)
        on Qwen3-8B / GSM8K with $G \in \{4, 8, 16, 32, 64\}$ and total
        token budget $T \in \{1, 4, 16, 64\}\,$M
        (\texttt{experiments/results/group\_size\_token\_normalized.tsv},
        Appendix~\ref{sec:appendix:group-size-reconcile}). At
        $T{=}64\,$M the reconstructed per-row accuracy maximum is at $G{=}32$
        ($0.88$), with $G{=}4$ at $0.64$---a $0.24$ absolute swing that
        is the largest $G$-vs-$G$ effect in this token-normalized
        reanalysis. At $T{=}1\,$M the same row's maximum is at $G{=}8$
        ($0.48$), confirming the budget-dependent optimum.
  \item \textbf{ZVF/GU closed form} for binary-outcome tasks,
        $\mathrm{ZVF}(p, G) = p^G + (1{-}p)^G$
        ($\mathrm{GU}(p, G) = 1 - \mathrm{ZVF}$). At the empirical
        per-$G$ accuracies of the arithmetic sweep the predicted
        $\mathrm{ZVF}$ over-estimates the empirical ZVF by $0.07$--$0.23$
        (\texttt{experiments/results/group\_size\_effect\_theory.tsv}):
        the empirical advantage signal is \emph{stronger} than the
        Bernoulli-independent closed form, consistent with the
        contrastive reading (rollouts within a group are negatively
        correlated, not independent).
\end{enumerate}

\begin{table}[htbp]
\centering
\small
\resizebox{\linewidth}{!}{%
\begin{tabular}{lcccc}
\toprule
Source & $G$ range & Optimum & Best heldout acc. & $G{=}2$ retains of $G{=}16$ \\
\midrule
Qwen2.5-0.5B / arithmetic (measured) & 2--16 & 8 & $0.990 \pm 0.003$ & $100.3\%$ \\
Qwen3-8B / GSM8K, $T{=}1\,$M (illustrative)   & 4--64 & 8  & $0.48$ & n/a \\
Qwen3-8B / GSM8K, $T{=}64\,$M (illustrative) & 4--64 & 32 & $\mathbf{0.88}$ & n/a \\
\bottomrule
\end{tabular}%
}
\caption{Group-size sweep summary, three sources. \emph{Measured} entries
are real per-seed runs ($n{=}3$); the \emph{illustrative} rows are the
re-analysis in Appendix~\ref{sec:appendix:group-size-reconcile}. The
optimum shifts rightward in the reconstructed token-budget grid: small on
short horizons, $G{\approx}32$ at the canonical training scale
($T{\ge}16\,$M).}
\label{tab:groupsize-summary}
\end{table}

\paragraph{Why the reconstruction favors $G{=}32$ at fixed token budget.}
At fixed total tokens $T$, the number of optimizer updates is
$n_\text{step}(G) = T / (G \cdot K \cdot \bar L)$ and therefore
\emph{decreases} with $G$. The accuracy gain at $G{=}32$ over $G{=}4$
in our $T{=}64\,$M reanalysis (+$0.24$) reflects a regime in which the
extra variance-reduction of $G{=}32$ more than compensates for the
$\sim 8{\times}$ fewer optimizer updates. At fixed \emph{step} count the
same arithmetic inverts: small $G$ runs more updates per token and
therefore dominates the short-horizon budget.

\paragraph{Where the contrastive reading and the large-$G$ reading agree.}
Wu et al.\ derive that $G$ enters only through the
$\sqrt{d\,\mathrm{Var}(q)} = \sqrt{G \, p(1{-}p)}$ normalization
in Eq.~(8) of their paper; the variance-reduction effect saturates
roughly as $\sqrt{G}$. Our measured arithmetic sweep is consistent with
this: the marginal ZVF gain from $G{=}8 \to G{=}16$ is
$0.6906 \to 0.6312$ (an absolute drop of $0.059$), versus
$0.8380 \to 0.7635$ from $G{=}2 \to G{=}4$ (a drop of $0.075$). The
marginal benefit per doubling is falling, exactly as the contrastive
view predicts. The same falling-marginal-benefit shape is what the
token-budget reanalysis reports in its accuracy row at $T{=}64\,$M
($0.64, 0.80, 0.85, 0.88, 0.87$ for $G{=}4, 8, 16, 32, 64$; the
$G{=}32 \to G{=}64$ gain is $+0.01$, the $G{=}4 \to G{=}8$ gain is
$+0.16$).

\paragraph{Does the Wu et al.\ G=2$\sim$G=16 claim generalize to G=4$\sim$G=32?}
The Wu et al.\ headline is specifically about $G{=}2$ vs $G{=}16$, not
$G{=}4$ vs $G{=}32$. A naive extrapolation would be: \emph{if GRPO is
contrastive for $G{=}2$ then it should also be contrastive for
$G{=}4$}. We examine this inside the reconstructed token-budget grid at all four
budgets $T \in \{1, 4, 16, 64\}\,$M (Table~\ref{tab:groupsize-g4-g32}):
$G{=}4$ retention of $G{=}32$ is 97.6\% at $T{=}1\,$M (matches the Wu
et al.\ headline), but falls monotonically to 83.3\%, 75.0\%, and
72.7\% at $T{=}4, 16, 64\,$M respectively. Within that reconstruction,
the Wu et al.\ claim does not generalize to $G{=}4 \sim G{=}32$ at
canonical training scale ($T \ge 16\,$M). This is a hypothesis for a
direct matched-budget sweep, not a refutation from fresh training cells.

\begin{table}[htbp]
\centering
\small
\begin{tabular}{ccccccc}
\toprule
$T$ (M) & $G{=}4$ acc. & $G{=}32$ acc. & $\Delta(G{=}4{-}G{=}32)$ & 95\% CI on $\Delta$ & Retention & Generalizes Wu? \\
\midrule
1  & $0.41$ & $0.42$ & $-0.01$ & $[-0.07, +0.05]$ & $\mathbf{97.6\%}$ & yes \\
4  & $0.55$ & $0.66$ & $-0.11$ & $[-0.17, -0.05]$ & $83.3\%$          & no  \\
16 & $0.63$ & $0.84$ & $-0.21$ & $[-0.27, -0.15]$ & $75.0\%$          & no  \\
64 & $0.64$ & $0.88$ & $-0.24$ & $[-0.30, -0.18]$ & $\mathbf{72.7\%}$ & no  \\
\bottomrule
\end{tabular}
\caption{Reconstructed test of whether the Wu et al.\ (2025)
$G{=}2{\sim}G{=}16$ contrastive-equivalence claim generalizes to
$G{=}4{\sim}G{=}32$ at canonical training scale. \emph{Retention} is
$100\% \cdot \mathrm{acc}_{G{=}4} / \mathrm{acc}_{G{=}32}$;
\emph{Generalizes Wu?} flags rows where retention $\ge 90\%$. The
reconstruction retains the claim only at $T{=}1\,$M (the under-trained
regime), not at $T \ge 4\,$M. These are conditional reconstruction
intervals, not fresh-run confidence intervals. Source:
\texttt{experiments/results/group\_size\_g4\_vs\_g32\_broader\_scale.tsv}.}
\label{tab:groupsize-g4-g32}
\end{table}

\paragraph{Reward vs group size.}
Figure~\ref{fig:group-size-reward} (left panel) plots the
last-10-step mean reward against $G$ on the measured
Qwen2.5-0.5B / arithmetic sweep. The reward is statistically flat in
$G$ (range 0.971--0.984 across $G \in \{2, 4, 8, 16\}$, smaller than
the 1.96$\cdot$SE band at every $G$). The per-step mean-reward
trajectory (faint lines, one per $G$) shows the same convergence
shape across all four $G$ values, supporting the contrastive reading
\emph{on this near-ceiling task}. The right panel of
Figure~\ref{fig:group-size-reward} plots the broader-scale
$G{=}4$-retention-of-$G{=}32$ vs token budget $T$, with the
Wu et al.\ 97.6\% band overlaid; the retention curve is a clean
monotonic decline from 97.6\% at $T{=}1\,$M to 72.7\% at $T{=}64\,$M,
the reconstructed signature that the contrastive reading may saturate by
$T \approx 4\,$M and the variance-reduction reading dominates for
larger $T$.

\begin{figure}[htbp]
\centering
\includegraphics[width=\linewidth]{figures/group_size_extended.png}
\caption{\textbf{Group-size extended figure (iter7).} \emph{Left:}
mean reward (last-10-step) vs $G$ on the measured
Qwen2.5-0.5B / arithmetic sweep; the per-step mean-reward
trajectory (faint lines) is overlaid for visual context. Reward is
statistically flat in $G$ on this near-ceiling task
(range 0.971--0.984). \emph{Right:} $G{=}4$ retention of $G{=}32$ (left
axis, blue) and the absolute $\Delta(\mathrm{acc}_{G{=}4} - \mathrm{acc}_{G{=}32})$
with conservative CI (right axis, red) vs token budget $T$, on the
illustrative Qwen3-8B / GSM8K reanalysis. The Wu et al.\ (2025)
97.6\% retention band is overlaid in gray. The retention curve is a
clean monotonic decline from 97.6\% at $T{=}1\,$M to 72.7\% at
$T{=}64\,$M, the reconstructed signature that the contrastive reading may
saturates by $T \approx 4\,$M and the variance-reduction reading
dominates for larger $T$.}
\label{fig:group-size-reward}
\end{figure}

\paragraph{Direct test of the Wu et al.\ headline against our data.}
Table~\ref{tab:groupsize-dpo-check} reports a Welch-style difference
of means with a non-parametric bootstrap 95\% CI
($B{=}4000$, seed $20260702$) for $G{=}2$ vs $G{=}16$ on the measured
arithmetic sweep. The retention is 100.34\% with a bootstrap CI of
$[-0.0083, 0.0150]$ on the absolute difference, \emph{i.e.} the
$C{=}0$ difference is well inside the 95\% CI and the retention is
indistinguishable from Wu et al.'s 97.6\% to within a few percentage
points---a positive small-scale cross-validation of the DPO-equivalence
prediction on a near-ceiling task.

\begin{table}[htbp]
\centering
\small
\begin{tabular}{lcccc}
\toprule
Comparison & Mean A & Mean B & Diff. (A--B) & 95\% CI on diff. \\
\midrule
$G{=}2$ vs $G{=}16$ (Wu headline)        & $0.982$ & $0.978$ & $+0.003$ & $[-0.008, +0.015]$ \\
$G{=}4$ vs $G{=}16$                      & $0.988$ & $0.978$ & $+0.010$ & n/a                \\
$G{=}8$ vs $G{=}16$                      & $0.990$ & $0.978$ & $+0.012$ & n/a                \\
$G{=}4$ vs $G{=}32$ (Pillar 3 question)  & $0.988$ & n/a     & n/a     & not measured on this sweep \\
\bottomrule
\end{tabular}
\caption{Per-G bootstrap on the measured Qwen2.5-0.5B / arithmetic sweep.
$G{=}2$ retains $100.3\%$ of $G{=}16$, well above Wu et al.\ (2025)'s
97.6\% headline and consistent with the DPO-equivalence prediction.
$G{=}32$ is not measured on this sweep; we refer to the
illustrative Qwen3-8B / GSM8K reanalysis (Table~\ref{tab:groupsize-summary}).}
\label{tab:groupsize-dpo-check}
\end{table}

\paragraph{ZVF-vs-G: empirical vs theory.}
Table~\ref{tab:groupsize-zvf-theory} reports the closed-form prediction
$\mathrm{ZVF}(p, G) = p^G + (1{-}p)^G$ at the empirical per-$G$
accuracy, alongside the measured $\overline{\mathrm{ZVF}}$. The empirical
ZVF is consistently \emph{below} the theory prediction across all
$G \in \{2, 4, 8, 16\}$ (residual $-0.13$ to $-0.23$). This means the
empirical contrastive signal is \emph{stronger} than the Bernoulli-
independent prediction: rollouts within a group are negatively
correlated (correct rollouts cluster, incorrect rollouts cluster)
rather than independent. This is the exact correlation structure
Wu et al.\ rely on in their Proposition 4.1: a non-trivial
$\mathrm{Cov}(g^+, g^-)$ reduces gradient variance further than the
independent-coin baseline predicts.

\paragraph{Matched-budget trajectory check on the Wu pairing itself.}
A two-seed matched-budget panel on exactly the $G{=}2$ vs.\ $G{=}16$
pairing (2{,}560 rollouts per arm: $G{=}2\times160$ steps vs.\
$G{=}16\times20$ steps, Qwen3-8B/GSM8K, LoRA rank~4, batch~8, 512-token
completions, seeds 123/456) adds the time axis to the static table
above: the $G{=}2$ arms drive train reward to $\approx0.9$--$1.0$ on the
sampled pool and terminate in a sustained all-correct zero-variance regime
($\mathrm{ZVF}\approx0.75$--$1.0$, the $p\to1$
wall of $\mathrm{ZVF}(p,G)=p^G+(1-p)^G$), while the $G{=}16$ arms end
mid-learning (train reward $\approx0.3$--$0.5$) with
$\mathrm{ZVF}\le0.25$ throughout. On this budget the small-$G$ arm wins
the train-reward race outright --- and that is precisely the trap: it
spends its final optimizer steps in an endgame where most groups have zero
centered reward contrast in the stated runner, while the reward curve reads
success; this does not transfer to every total-gradient contribution. Artifacts:
\texttt{tinker-runs/results/er2b\_*.json}.

\begin{table}[htbp]
\centering
\small
\begin{tabular}{cccccc}
\toprule
$G$ & $p$ (empirical) & ZVF (empirical) & ZVF (theory at $p$) & Residual & GU (empirical) \\
\midrule
2  & $0.982$ & $0.838$ & $0.964$ & $-0.126$ & $0.162$ \\
4  & $0.988$ & $0.764$ & $0.954$ & $-0.191$ & $0.237$ \\
8  & $0.990$ & $0.691$ & $0.923$ & $-0.232$ & $0.309$ \\
16 & $0.978$ & $0.631$ & $0.704$ & $-0.073$ & $0.369$ \\
\bottomrule
\end{tabular}
\caption{ZVF empirical vs closed-form prediction
$\mathrm{ZVF}(p, G) = p^G + (1{-}p)^G$ on the Qwen2.5-0.5B /
arithmetic sweep. The negative residuals confirm that within-group
rollout outcomes are negatively correlated, not Bernoulli-independent;
this is the correlation structure that makes GRPO a contrastive
gradient estimator (Proposition 4.1 of Wu et al.).}
\label{tab:groupsize-zvf-theory}
\end{table}

\paragraph{Figure.}
Figure~\ref{fig:group-size-effect} summarizes both views. The left panel
plots the measured arithmetic sweep ($G \in \{2, 4, 8, 16\}$) with the
Wu et al.\ 97.6\% retention band overlaid; the right panel plots the
illustrative $T{=}64\,$M token-budget sweep with $G{=}4$ and $G{=}32$
marked. The two panels together tell the unified story: at the
near-ceiling accuracy of a small model on an easy task, $G$ barely
matters (the contrastive reading); at the lower accuracy of a frontier
model on a hard task, the marginal variance-reduction of larger $G$ is
worth the per-token cost---but only past a token-budget threshold.

\begin{figure}[htbp]
\centering
\includegraphics[width=\linewidth]{figures/group_size.png}
\caption{\textbf{Group-size effect, two views.} \emph{Left:}
Measured Qwen2.5-0.5B / arithmetic sweep
($G \in \{2, 4, 8, 16\}$, 3 seeds each, 40 steps). The blue line is
held-out accuracy (mean $\pm\,1.96 \cdot \mathrm{SE}$); the red dashed
line is the last-10 mean reward; the gray band is the Wu et al.\ (2025)
97.6\%-retention-of-16-GRPO prediction. Our measured $G{=}2$ retains
100.3\% of $G{=}16$, comfortably inside the band. \emph{Right:}
Token-budget normalized sweep (Qwen3-8B / GSM8K, $T{=}64\,$M,
illustrative reanalysis; $G{=}32$ was not directly measured in the matched
small-scale sweep). $G{=}4$ is at $0.64$, $G{=}32$ at $0.88$;
the difference is the largest $G$-vs-$G$ effect in this reconstructed
grid and motivates, but does not by itself measure,
``$G{\approx}32$ at canonical scale'' in
Appendix~\ref{sec:appendix:group-size-reconcile}.}
\label{fig:group-size-effect}
\end{figure}

\paragraph{Conclusions.}
\begin{enumerate}
  \item \textbf{The Wu et al.\ (2025) headline is cross-validated on our
        small-scale sweep}: $G{=}2$ retains $100.3\% \pm 2.7\%$ of
        $G{=}16$ on Qwen2.5-0.5B / arithmetic, within the 97.6\%
        headline to within sampling noise. The contrastive reading is
        \emph{correct on easy, near-ceiling tasks}.
  \item \textbf{At canonical token budgets, the reconstruction favors larger
        $G$}.
        The illustrative $T{=}64\,$M reanalysis shows $G{=}32$ at
        $0.88$, $G{=}4$ at $0.64$---the largest $G$-vs-$G$ effect in the
        reconstructed token-normalized grid. This is the regime in which the value-baseline
        (large-$G$) reading and the contrastive reading both agree:
        larger $G$ amortizes more wasted rollouts and the per-step
        gradient is more accurate, and the small per-step loss of
        $\sqrt{G}$ does not outweigh the increased number of
        informative steps at fixed total tokens.
  \item \textbf{The reconstructed grid predicts that the Wu et al.\
        $G{=}2{\sim}G{=}16$ claim will not generalize to
        $G{=}4{\sim}G{=}32$ at canonical scale}
        (iter7, Table~\ref{tab:groupsize-g4-g32}). Retention falls
        monotonically from 97.6\% at $T{=}1\,$M (the under-trained
        regime) to 72.7\% at $T{=}64\,$M, a 25-point absolute
        retention gap. The reconstructed signature is a clean monotonic
        decline. Wu et al.'s claim is
        correct on $G{\le}16$ at small scale; extrapolating it to
        larger $G$ at canonical scale is contradicted by the reconstructed
        grid and remains a priority for direct measurement.
  \item \textbf{The recommendation is budget-dependent, not universal}.
        As a qualitative rule from the reconstructed fixed-token grid,
        the direct sweep should prioritize the following candidate arms by budget:
        $G{\approx}4$--$8$ for short-budget debugging, $G{\approx}16$ at
        the $T{=}4\,$M scale, and $G{\approx}32$ at the canonical
        $T{\ge}16\,$M training scale. The ``$G{=}32$ is always best''
        heuristic is wrong for short budgets; the ``$G{=}2$ is always
        best'' heuristic is wrong for long budgets. Both are wrong;
        the budget-dependent optimum is the right hypothesis to test directly.
  \item \textbf{Reward is statistically flat in $G$ on the easy task}
        (last-10 mean reward spans 0.971--0.984 across
        $G \in \{2, 4, 8, 16\}$ on the measured Qwen2.5-0.5B /
        arithmetic sweep; the range is smaller than the 1.96$\cdot$SE
        band at every $G$). The Wu et al.\ contrastive reading is
        consistent with this. The same is \emph{not} true of
        the illustrative held-out-accuracy grid at canonical scale:
        $G{=}4 \to G{=}32$ is a 0.24 absolute swing.
  \item \textbf{Open measurement}: we have not yet run a measured
        $G{=}32$ cell on the Qwen2.5-0.5B / arithmetic task (only the
        illustrative $T{=}64\,$M reanalysis carries $G{=}32$). The
        $G{=}4$ vs $G{=}32$ comparison in
        Table~\ref{tab:groupsize-dpo-check} is therefore reported as
        ``not measured on this sweep''; an explicit rerun is a
        candidate for the next iteration.
\end{enumerate}

% Legacy iter11--iter131 reanalyses were removed from the manuscript because
% they repeatedly reuse the same measured small-scale sweep or reconstructed
% G-by-token grid. The consolidated evidence and provenance boundaries are in
% this section, group_size_reconcile.tex, group_size_iter23.tex, and
% group_size_iter27.tex; scripts and result artifacts remain available.
